% Vietnamese translation for OI Public Library
% Source-File: IOI中国国家候选队论文/2009/张昆玮/数学归纳法与解题之道.ppt
\documentclass[11pt]{article}
\usepackage{oipl-vn}

\title{Quy nạp toán học và con đường giải bài\\{\large Ghi chú theo slide}}
\author{Zhang Kunwei\\Trường Thực nghiệm Sơn Tây\\Huấn luyện viên: Tang Wenbin, Hu Weidong}
\date{16 tháng 12 năm 2008}

\begin{document}
\maketitle

\origtitle{Mathematical induction and problem solving}
\sourcepath{IOI-China-Candidate-Team-Papers/2009/Zhang-Kunwei/induction-slides.ppt}

\section{Vai trò của quy nạp}

Bài trình bày đặt quy nạp toán học vào bối cảnh giải bài thi tin học: không
chỉ để chứng minh mệnh đề đã biết, mà còn để tìm cấu trúc, thiết kế thuật
toán xây dựng và kiểm tra vì sao một lựa chọn cục bộ là hợp lệ.

\section{Các cách dùng trong thuật toán}

\begin{enumerate}
  \item Chứng minh tính đúng đắn của tham lam: sau khi xử lý một phần nhỏ,
  chứng minh có thể mở rộng thành lời giải tối ưu cho phần lớn hơn.
  \item Xây dựng cấu trúc: từ trường hợp \(n-1\), chỉ ra cách thêm phần tử
  thứ \(n\) mà vẫn giữ bất biến.
  \item Tối ưu hóa: chọn đối tượng quy nạp phù hợp hơn, đổi hướng quy nạp,
  hoặc tăng cường giả thiết để bước chuyển trở nên đơn giản.
  \item Phát hiện giới hạn: có những bài phụ thuộc lịch sử hoặc có vòng lặp
  trạng thái, nên quy nạp trực tiếp không còn là công cụ thuận tiện.
\end{enumerate}

\section{Gợi ý đọc slide}

Khi gặp một thuật toán trong slide, nên hỏi ba câu: cơ sở quy nạp nằm ở đâu,
bất biến nào được giữ sau mỗi bước, và giả thiết quy nạp có đủ mạnh để chứng
minh bước chuyển hay không. Nếu câu trả lời thứ ba là không, thường cần đổi
đại lượng quy nạp hoặc thêm thông tin phụ vào trạng thái.

\end{document}
